\addbibresource{references.bib}

\title{Reflections on Capital Punishment}
\author{Jason Reeves (October 2026)}
\date{}

\begin{document}

\ifdefined\embedversion
  \pagestyle{empty}
\fi

\maketitle
\extrafloats{100}

\begin{abstract}
Christians are taught not to be ashamed of the gospel\footnote{For I am not ashamed of the gospel, for it is the power of God for salvation to everyone who believes, to the Jew first and also to the Greek.\scriptref{Romans 1:16 \biblever{nasb}}}, but many who heed this command fail to realize that included in this command is the strong implication that we also not be ashamed of the law\footnote{Psalm 19, Psalm 119:6, 119:46, 119:80, 119:97}, in which the gospel is rooted.  The Old Testament moral law is binding on us as Christians today, as well as the implications and precepts of the obsolete legal codes of the Israelite theocracy and fulfilled ceremonial laws\footnote{This is known in Covenant Theology as the \textit{tripartite division of the law}.  The law is divided into the moral law (perpetual and binding), the civil laws of the theocratic Israelite nation (expired with that nation), and the ceremonial and sacrificial laws (fulfilled in Christ).  Though only the moral law is directly binding, principles and precepts gleaned from the inoperative laws still reflect the character of God, Who does not change.}.  Christians are not \textit{required} to obey and uphold these laws as a \textit{condition} of salvation, but God's word teaches us that as we grow in faith and are further sanctified, our lives will show the fruit of obedience\footnote{Galatians 5:16-26}.  We would do well to study and apply God's law to our lives\footnote{Psalm 119:104}.

Capital punishment, or the death penalty, is a practice sanctioned by the Old Testament.  Paul reaffirms the right of the civil magistrate to carry out the death penalty, and comments that the magistrate is ordained by God so to do.  Christians, who are called to love the things that God loves and hate the things that God hates, must uphold and affirm this right and duty of the civil magistrate to carry out such punishments.  We are called to rejoice when God's justice is carried out.  Yet, though we must rejoice for this purpose, we must balance this joy with sorrow for the condemned, who bears God's image.
\end{abstract}

\newthought{As Francis Schaeffer said, "How should we then live?"} How do we uphold the principle of capital punishment while avoiding the danger of creating a violent, bloodthirsty, vengeful society, as well as cultivating these traits in ourselves and encouraging or ignoring these traits in fellow believers?  Clearly this is neither the type of society that God would have His people to create, nor the type of people He has desired to call out for Himself -- these aspirations are not worthy of those for whom Christ died.  God's word speaks against wanton violence, bloodthirsty attitudes, and personal vengeance.  Boasting, and by extension gloating over a defeated enemy is also spoken against.  If we are to envision a society in which God would be pleased, such a society must be one that is, among other things: merciful, while not neglecting justice; humble, yet boasting in the Lord and insisting on His ways; and peaceful, yet willing to use violence when necessary.  These characteristics must also be those for which God's blood-bought children strive, and for which the church must encourage and equip them.

\ifdefined\embedversion\newpage\fi
\section{What Changed?}
There was a time when public, violent spectacles were simply accepted as normal.  Jesus Christ entered the world during this time, and God's chosen means for the death of Christ was an official execution accompanied by a public spectacle.  Christ's followers fared no better in the centuries closest to His death, burial, and resurrection.  A child's understanding of church history should necessarily include stories of the Roman Colosseum and Circus Maximus, where Christians were publicly executed using the most gory methods conceivable, after having been prosecuted for what the Romans considered to be the crime of atheism, by which they meant that the Christians refused to worship the established Roman pantheon\footnote{There were other charges levied, such as \textit{odium generis humani}, i.e. "hatred of the human race", according to Tacitus.  Additionally, charges of not being willing to sacrifice to Caesar or accord him as a god, sacrilege, blasphemy, etc.}.

In history it can seem like nothing happens for centuries, then everything happens all at once.  In the midst of and subsequent to the American Great Awakenings, Western Civilization experienced many reforms fueled by Christian sentiment.  Many of these reforms were good.  Some were not. William Wilberforce and John Newton fought against the transatlantic slave trade (and by extension slavery in general) from a Christian perspective and won.  A bloody war was fought in America over slavery, and as a result of this war, the so-called \textit{peculiar institution} was abolished, and good riddance.  Thomas Sowell pointed out that the institution of slavery, which had been unquestioned for millennia, had been eradicated within the space of about a century\footcite{sowellblack}.  Robert Fogel, who conducted research on American slavery and was himself an atheistic Jew, concluded that it was Christianity that actually accomplished this monumental civilizational achievement\footcite{carpentersecularjew}\footcite{carpenterapologist}. 

So we see that it is not always the case that "the old ways are the good ways".  To be \textit{conservative} one must be attempting to \textit{conserve} something.  What one proposes to conserve can make a conservative of one era completely distinct from a conservative of another era.  The Pharisees were conservatives during the time of Christ.  They were those that insisted on the old ways of Judaism and temple worship, continuing this even after the Christ had come.  In the case of the Pharisees, the old ways were definitely \textit{not} the good ways.  Similarly, those clinging to slavery at a time when God was making it clear to Christians worldwide that the practice had to end were in the wrong.  We may be able to \textit{understand} their position, but we cannot \textit{defend} and we should never consider \textit{emulating} it.  

But then consider Machen's monumental work, "Christianity And Liberalism"\footcite{machenchristianity} in which he casts early 20th century liberalism as a religion in its own right distinct from Biblical Christianity.  Machen is a conservative, seeking to conserve concepts of biblical inerrancy and sufficiency over against the forces assaulting these ideas due to Enlightenment ideals run amok, which resulted in so-called "higher criticism".  In this case, the old ways \textit{were} the good ways.  In fact, I believe that in most situations throughout history, asserting the old ways is \textit{generally} the right move.  But not always.  Sometimes, merely asserting the old ways without consideration for why they were abandoned or what the Bible says is a path to grievous error.

I believe that as Western sensibilities have been sanctified due to the leaven of the Christian gospel, many practices antithetical to Christian philosophy that had existed for some time disappeared.  The Enlightenment, though atheistic at its core, owes its rigid epistemic standards and sense of truth to numerous Christian scientists who were seeking order in creation because they \textit{expected} to see order -- because they read in God's word that He created an orderly universe. That Christianity is based on the bedrock truth of Scripture, Western civilization is based on Christianity, and that the Enlightenment was based on the ravings of God-hating French philosophes are good points to bear in mind.  The Enlightenment, though resting on totally separate foundations, benefited from and twisted Christian precepts, even as it denied God's very existence.  We should not be surprised when, at times, Christian and Enlightenment thought intersect.  That both slavery and spectacular public executions have been spoken against by both Christian and Enlightenment reasoning doesn't make either a "Christian" or "Enlightenment" issue primarily.  Let us not forget the Enlightenment gave us a new "scientific" execution instrument, the guillotine, which was used to great effect at many execution spectacles.  Let us also reflect on how Christian communities have been speaking out against public brutality, including slavery, for some time before the Enlightenment. Indeed, these \textit{are} Christian issues, and if the Enlightenment has any claim on them, it is at best derivative\footnote{I beg the reader's pardon for this digression into the nature of conservatism and the Enlightenment.  What I am attempting to head off is the accusation that a particular reform was not a result of "gospel leavening" but rather of "Enlightenment liberalism".  To unpuzzle this conundrum for any reform, one only needs to determine if the reform is supportable from Scripture and examine what the general thoughts of God's people have been through the centuries.  No doubt there are "false sons" in the pale of the church, as the hymn says, so one must be judicious.}.

According to the historical record the first New England sermon to openly question the practice of public executions was delivered in Massachusetts in 1801 -- on the occasion of the hanging of Jason Fairbanks\footcite[The main thrust of this dissertation is that the \textit{liberalizing} of Christianity in New England led to the shift away from executions being open to the general public. I believe that conservatives \textit{and} liberals were able to unite on getting rid of the \textit{spectacle} aspect.  The liberals pressed forward from this and worked effectively to abolish the death penalty entirely where it was in their power so to do.]{belczakchanging}.  Up to this point in America, as elsewhere in Christendom, public executions were treated as solemn events.  There would be an execution sermon which would emphasize God's sovereignty and the duty of the magistrate to wield the sword.  Those attending would be led in singing hymns and psalms.  It was almost like a church service.  Unfortunately, there were increasingly more cases such as the Fairbanks hanging where thousands of people who had no connection to the crime or the place showed up merely to watch a hanging for sport, entertainment, or pleasure. Belczak remarks:
\begin{quote}
The execution ritual no longer fulfilled the culturally prescribed role that it once had, or at least had aspired toward.  What was supposed to be a religious or civic ritual had become a kind of carnival; the thousands came to Dedham as gawking onlookers rather than pious observers on that September day in 1801.  This reality was not lost on [Thacher], who, facing this discrepancy between the ideal and the reality surrounding the execution ritual, argued for the removal of this corrupting sight away from the eyes of the social body.  [Thacher] even implored the audience to examine the situation themselves and ask "whether any considerate man present the last week, at the execution, can affirm, that any religious or moral benefit was to be gained from what he saw or heard."\footcite{belczakchanging}
\end{quote} 

Thomas Thacher, though a liberal proto-Unitarian, never voiced opposition to the death penalty itself, only to \textit{how} the penalty was carried out, with concerns about the potential negative effects on a society which sees such events not as a solemn carrying out of God's will by the magistrate, but as a means of casual entertainment or perhaps even vicariously experienced retribution.  Jonathan Going, who, according to Belczak, delivered the last known execution sermon in New England, was a conservative Reformed Baptist.  This sermon was preached in 1825, also in Massachusetts, at the execution of Horace Carter.  Going suggested that it might be better for public morals if executions were held behind prison walls with only the necessary people present\footcite[This is largely my position.  The death penalty should be carried out.  Victim families, law enforcement, members of the press should be able to attend. Nothing should be transmitted live, because anything transmitted live can be recorded and replayed.  No video or audio should be taken.  Nothing that could give any kind of relevance or martyrdom to the condemned prisoner should be allowed.]{belczakchanging}.

In England, conservatives such as the aforementioned William Wilberforce fought for reform of the so-called "Bloody Code" which at its height defined more than 200 capital crimes.  Samuel Wilberforce, William's son, said in the UK House of Lords in 1854:
\begin{quote}
Such scenes as lately occurred at executions at Monaghan and other places had forced a conviction on the English mind that people were not deterred from crime by such spectacles, and that, therefore, the continuance of capital punishment, as at present conducted, would have the effect of encouraging the morbid sentimentality in favour of the criminal, and to lead men to doubt whether the highest duty committed to man—that of wielding the sword of justice —was committed by God to man for civil government, and to doubt whether that civil government had not the right to take away life, not only to deter men from crime, but as an execution of God's sentence against great offenders. Nothing could be more injurious or more enervating than to put down capital punishments in consequence of the morbid sentimentality to which he had referred. A few more such scenes as the noble and learned Lord had referred to would have the effect of making men's minds recoil from that which he (the Bishop of Oxford) believed to be essential to the highest principles of justice, as well as to the necessities of human expediency\footcite{wilberforcecapitalexecutions}.
\end{quote}

After the 1800s, there was a decline of the practice of public executions across the United States and the world.  I attribute this primarily to the influence of the gospel on American society, and the prominence and power of America in the world.  The image of the gospel leaven\footnote{Matthew 13:33, Luke 13:20-21} is not one of instantaneous and simultaneous effect of the gospel everywhere, but rather the image is of a slow, yet inexorable, and also inevitable progress of the sanctifying effects of the gospel according to God's good pleasure.  This sanctification does not \textit{produce} salvation, but rather establishes an environment which is favorable toward salvation.  Why does every civilized nation in the world have murder statutes and prosecute people for murder? Even in nations which do not acknowledge Christ, such statutes exist. Although God's common grace might have enlightened totally pagan cultures to create such statutes, yet their wide prominence, proliferation, and maintenence has been due in large part, I believe, to this effect of gospel leaven by way of the influence of Western civilization.  Additionally, cultures which may have been forced to live for a time under such laws have recognized the goodness and general equity of them. 

But this gospel influence is not constant.  For reasons known only to God, it sometimes advances, and sometimes retreats.  We are only promised that it will leaven the whole loaf in God's timing.  

In early 19th century New England, I believe the clergy and the magistrate, at least, had observed the potentially harmful effects a public execution could have on society.  The clergy called for reforms, the magistrate obliged, and the masses lagged somewhat behind.  After living without public executions for a few generations, the masses no longer desired them.

\ifdefined\embedversion\newpage\fi
\section{What Does Scripture Say?}

\subsection{The Biblical Warrant}
There are some crimes which deserve the death penalty.  The Old and New Testaments provide mechanisms for communities to put members to death for especially heinous crimes.  In the Old Testament, this was done within the covenant community of Israel, with the witnesses throwing the first stone, and the entire congregation participating afterwards\footnote{The hand of the witnesses shall be first against him to put him to death, and afterward the hand of all the people.  So shall you purge the evil from your midst.\scriptref{Deuteronomy 17:7 \biblever{nasb}}}.  In the New Testament, the civil magistrate is the vehicle ordained by God to perform this function\footnote{\ldots\ for it [the magistrate] is a minister of God to you for good.  But if you do what is evil, be afraid; for it does not bear the sword for nothing; for it is a minister of God, an avenger who brings wrath on the one who practices evil.\scriptref{Romans 13:4 \biblever{nasb}}}.  We participate in such penalties by consenting to live under the laws of and by paying taxes to the magistrate.

The broad categories of crimes assigned the death penalty included: offenses against God (e.g. Exodus 22:20), various sexual sins (e.g. Leviticus 20:10), murder (e.g. Exodus 21:12), kidnapping (e.g. Exodus 21:16), and rebellion against God-ordained authority (e.g. Deuteronomy 17:12).

The Bible also anticipates that capital cases will be taken care of in an expedient manner.  For better or for worse, modern governments do not operate in this way.  We give a condemned man every possible opportunity for appeal.  Perhaps it is good to be somewhat circumspect and make sure that the one convicted is actually guilty.  Yet Ecclesiastes 8:11 says, "Because the sentence against an evil deed is not executed quickly, therefore the hearts of the sons of men among them are given fully to do evil." Certainly, the deterrent function of capital punishment, used in Scripture to recommend the practice\footnote{Deuteronomy 13:11, 17:12-13, 19:19-20, 21:21}, is significantly diminished when that punishment is not quickly carried out.  

\subsection{Justice, Not Vengeance}
The death penalty properly carried out is never characterized in Scripture as feeding some perceived need for personal or corporate retribution or vengeance.  God alone is responsible for wielding the sword of vengeance\footnote{Never pay back evil for evil to anyone.  Respect what is right in the sight of all men.  If possible, so far as it depends on you, be at peace with all men.  Never take your own revenge, beloved, but leave room for the wrath of God, for it is written, `Vengeance is Mine, I will repay,' says the Lord.\scriptref{Romans 12:17-19 \biblever{nasb}}}.  He is the only one who can righteously execute justice in this way.  The death penalty is characterized, rather, as the solemn carrying out of God's justice.  Any personal or corporate retributive spirit taints and blunts the just nature of the penalty and potentially brings God's just condemnation on such motives.  We are never as a society going to get this completely right. This inability to attain perfection should not deter us from pursuing obedience, either in civil or personal life.  

\subsection{Mourning And Rejoicing}
Scripture teaches us not to rejoice at the fall of the wicked\footnote{Do not rejoice when your enemy falls, and do not let your heart be glad when he stumbles; or the \textsc{Lord} will see it and be displeased, and turn His anger away from him.\scriptref{Proverbs 24:17-18 \biblever{nasb}}}, but also teaches us that the results of the wicked falling are good and worthy of praising God\footnote{Rejoice over her, O heaven, and you saints and apostles and prophets, because God has pronounced judgment for you against her.\scriptref{Revelation 18:20 \biblever{nasb}}}\footnote{The righteous will rejoice when he sees the vengeance [of God]; he will wash his feet in the blood of the wicked.  And men will say, `Surely there is a reward for the righteous; surely there is a God who judges on earth!'\scriptref{Psalm 58:10-11 \biblever{nasb}}}.   In the case of the execution of a murderer, we must learn to mourn the loss of a fellow image bearer\footnote{Say to them, `As I live!' declares the Lord \textsc{God}, `I take no pleasure in the death of the wicked, but rather that the wicked turn from his way and live.  Turn back, turn back from your evil ways!  Why then will you die, O house of Israel?'\scriptref{Ezekiel 33:11 \biblever{nasb}}} while simultaneously celebrating the upholding of God's justice. These two concepts can be confused, and people might come away believing that they need to do one to the exclusion of the other. However, these two Biblically-prescribed reactions are distinct.  Yet, they must be kept in balance.

Why should we mourn the loss of a murderer or a terrorist? First, it is tragic whenever a human, with all the potential God has placed within him, chooses to use that potential for evil.  That is a shame.  That is a missed opportunity.  That is an affront to God.  What a waste.  Second, we are called to have empathy.  In recent years, we've seen toxic empathy as well as suicidal empathy -- empathy that is simply performative and self-serving, as well as empathy that is completely destructive to all involved -- neither of which actually provides help to the one suffering.  But Christian empathy is a real thing which we are commanded to exhibit, and it is intended to bless others and honor God\footnote{Bless those who persecute you; bless and do not curse.  Rejoice with those who rejoice, and weep with those who weep.  Be of the same mind toward one another; do not be haughty in mind, but associate with the lowly.  Do not be wise in your own estimation.\scriptref{Romans 12:14-16 \biblever{nasb}}}.  Finally, if we are tempted to exhibit our own personal wrath for the condemned, meditating on their humanity and manifesting empathy for the condemned could help to offset such impulses.  

\subsection{Bloodlust And Conscience}
Efforts to create a just and merciful society must necessarily exclude the cultivation of bloodlust, the inuring of consciences to casual brutality, and the encouragement of vengeful impulses which are antithetical to the Christian mindset.  If we are to create communities which operate on Christian principles, then official acts must take into consideration how such acts might affect the consciences and behaviors of the populace.  The death penalty is supposed to act as a deterrent, but how that penalty is carried out can encourage or discourage Christian impulses.  The magistrate should be encouraged to choose prudently.

God requires Christians to work toward a just society, but a society characterized by justice must not neglect mercy.  A wholly just yet unmerciful society would be as much of a hell on earth as one that was wholly merciful yet unjust.  James warns against withholding mercy.  He says that those who fail to show mercy will have such mercy withheld from them\footnote{So speak and act as those who are to be judged by the law of liberty.  For judgment will be merciless to one who has shown no mercy; mercy triumphs over judgment.\scriptref{James 2:12-13 \biblever{nasb}}}.  Similarly, Hosea relates the principle of bloodshed begetting bloodshed\footnote{There is swearing, deception, murder, stealing, and adultery.  They employ violence, so that bloodshed follows bloodshed.\scriptref{Hosea 4:2 \biblever{nasb}}}. Ezekiel quotes God himself expressing His expectation that His people should hate bloodshed. Because His people did not hate bloodshed, God said, "therefore bloodshed will pursue you."\footnote{\ldots\ therefore as I live,' declares the Lord \textsc{God}, `I will give you over to bloodshed, and bloodshed will pursue you; since you have not hated bloodshed, therefore bloodshed will pursue you.'\scriptref{Ezekiel 35:6 \biblever{nasb}}}

It should be noted that this principle has been overstated in our modern society.  Many of those who say "violence begets violence" rail against faithful Christian parents who use the rod to discipline their children, for example.  In matters of war, they are pacifists.  In matters of self-defense or defense of the weak and innocent, they prefer to allow evil to triumph so that good may result.  With Paul, I say, "Their condemnation is deserved." Many of those same people would also be against the death penalty in any instance.  I believe that what the Bible is speaking against is \textit{wanton} or \textit{purposeless} violence, violence without meaning or significance, or perhaps with a meaning or significance not sanctioned by God.  Parents disciplining their children with the rod is a form of violence, couched in love when properly carried out, with means and ends well supported by Scripture.  Yet, there is a \textit{wrong} way to approach disciplining one's children.  There is an \textit{unrighteous} way to carry out such discipline, even though the Bible prescribes the rod to drive rebellion far from the heart of the child.

So it is with capital punishment.  This is a form of violence, prescribed by God in the Bible as a means of executing His justice and deterring future crime.  Yet, there is an unrighteous way to carry out executions.  There is an unrighteous way for Christians to think and talk about executions.  

With regard to the execution of Nidal Hasan in 2026, there was debate about whether or not the execution should have been streamed live for anyone around the world to see.  Two justifications \textit{why} it should be streamed live were common in right-wing spaces on the Internet: 1) "we had to watch a live execution on 9/11!" and 2) "we had to watch a live execution when Charlie Kirk was killed!"  I use an exclamation point after each of these to indicate that there was usually a high emotional state behind these declarations.  While emotion can be an appropriate and effective communicator, it can also cloud the reasoning, as it seems to have done here.  An execution, or a carrying out of a properly adjudicated death sentence, is a \textit{good} thing, and it is done by society, collectively.  It is a legitimate use of force.  It is a practice that is ordained by God.  To call the murders which happened on 9/11 or the murder of Charlie Kirk "executions" has the effect of legitimizing these murders.  God sanctions the death penalty.  He does not sanction murder.  Additionally, attempting to justify making Hasan's execution available to anyone worldwide based on "they do it, too!" is a textbook employment of the \textit{tu quoque} logical fallacy.  \textit{Tu quoque} is Latin for "you, too" and is a subspecies of \textit{ad hominem} attack.  When asking the question "Is this right?" -- the justification of somebody else having done it or not is illogical and inadmissible\footnote{Unless, of course, the \textit{hominem} in question is Christ or another positive Biblical example.}.  As Christians, how would God have us live?\footnote{There are other reasons I believe this is a bad idea, such as the fact that Hasan's execution footage could be used in many unsavory ways in which it was not intended (promotion of martyrdom, fomenting unrest, radicalization efforts, etc.).  I'm thankful the decision was made \textit{not} to stream the event live.}

Proverbs lists being willing to "hasten to shed blood" as one of the traits of the wicked\footnote{My son, do not walk in the way with them [the wicked].  Keep your feet from their path, for their feet run to evil and they hasten to shed blood.\scriptref{Proverbs 1:15-16 \biblever{nasb}}}.  We should be very careful as a society and as individuals not to simply react in anger when a crime is committed. Only after due process, calm consideration, and giving the issue the solemnity it deserves should we consider the punishment, how it may be righteously carried out, and how we as Christians should think, feel, speak, and act about it\footnote{This you know, my beloved brethren. But everyone must be quick to hear, slow to speak and slow to anger; for the anger of man does not achieve the righteousness of God. Therefore, putting aside all filthiness and all that remains of wickedness, in humility receive the word implanted, which is able to save your souls. But prove yourselves doers of the word, and not merely hearers who delude themselves.\scriptref{James 1:19-22 \biblever{nasb}}}.

\ifdefined\embedversion\newpage\fi
\section{Conclusion}
The death penalty should be heart-wrenching, even when the condemned is a murderer or a terrorist.  Individuals and nations will be held to account for the diligence they employ to get this right, or else for the flippancy and disregard demonstrated.  Consequently, they will also be judged for their \textit{laxity} in enacting just punishment for crimes or being slow to carry out such justice.  If life, even the life of a condemned criminal, is approached carelessly, or worse, brutally, it cheapens all life.  Remember Christ had mercy on a condemned criminal, and though that criminal was guilty and Christ did not erase his physical penalty, yet He afforded the man dignity and compassion. 

As Christians, we should support the death penalty for heinous crimes while preserving dignity and compassion for the condemned, knowing that we will make mistakes from time to time, yet laboring toward a clear conscience fueled by due diligence.  In order that the God-ordained deterrent effect of capital punishment not be blunted, it should be carried out quickly, yet preserving due process.  Executions should happen in a private environment using methods of death which are certain in their action yet preserve the dignity of all -- the condemned, law enforcement, etc.  Press, law enforcement, public officials, and victim families should be allowed to witness.  No recordings should be allowed, lest the condemned attain undue significance in the public eye.  In order to encourage the best Christian impulses in society, the magistrate must consider the consciences of that society.

\end{document}
